\documentclass[10pt,a4paper]{amsart}
\usepackage[margin=19mm]{geometry}
\usepackage{amsmath,amssymb,mathtools}
\usepackage[scr=rsfs,cal=euler]{mathalfa}
\usepackage{xfrac}
\usepackage[unicode,hidelinks,bookmarks=false]{hyperref}
\newcommand{\ie}{i.e.\ }
\newcommand{\cf}{cf.\ }
\newcommand{\resp}{respectively\ }
\newcommand{\Subsection}{\subsection}
\providecommand{\nobreakdash}{\nobreak\hbox{-}}
\setlength{\emergencystretch}{2em}
\multlinegap0pt
\usepackage{bm}
\usepackage{bbm}
\usepackage{dsfont}
\usepackage{xspace}
\usepackage{cancel}
\usepackage{mathtools}
\usepackage{amsthm}
\makeatletter
\@ifundefined{theorem}{\newtheorem{theorem}{Theorem}}{}
\@ifundefined{proposition}{\newtheorem{proposition}[theorem]{Proposition}}{}
\makeatother
\newcommand{\C}{\mathbb{C}}
\newcommand{\R}{\mathbb{R}}
\newcommand{\Z}{\mathbb{Z}}
\title[Conjugate phase retrieval on graphs and with applications in]{Conjugate phase retrieval on graphs and with applications in shift-invariant spaces}
\author{Cheng Cheng, Baixiang Wu, Jun Xian}
\date{Source: arXiv:2507.22468v1 (CC BY 4.0)}
\begin{document}
\maketitle

\noindent\emph{Source and licence.} Conjugate phase retrieval on graphs and with applications in shift-invariant spaces. Version arXiv:2507.22468v1 (CC BY 4.0), distributed under the Creative Commons Attribution 4.0 International licence, as recorded in the arXiv licence metadata for this exact version and shown on the abstract page https://arxiv.org/abs/2507.22468v1. Full rights notice: licenses/arxiv-2507.22468-CC-BY-4.0.txt.\par\smallskip

\noindent\textit{Editorial scope.} Counted unit: the Theorem whose source label is \texttt{cpr.adj.thm} in the article (source0.tex lines 785--787, source label \texttt{cpr.adj.thm}), delivered together with the article's own complete original proof of it (source0.tex lines 788--863, one \texttt{proof} environment of 76 lines). No mathematics is written, repaired or re-derived: statement and proof are the article's own text line for line. Required context retained uncounted: sampling.def.eq1 (lines 142--144); prop2 (lines 663--665); eq91 (lines 776--778). Every definition, display and label of those lines is retained unchanged. Omitted from this package and not redistributed: 1--141, 145--662, 666--775, 779--784, 864--1417. Nothing inside a retained excerpt is omitted or altered: every retained body is delivered exactly as the article's lines are written, so no omission receipt of any kind accompanies this package. Citations: the retained excerpts carry no citation command at all, so no bibliography entry of the article is transcribed and no reference is redistributed. Reference adaptation: none was needed and none was made; every reference inside the delivered unit resolves inside it, so no unresolved reference or citation is produced. Printed slips kept literal: no printed word, symbol, hypothesis or number of the counted statement or of its proof was altered. Layout and class: the article's own class is \texttt{amsart} with packages including geometry, graphicx, float, subfigure, hyperref, multirow, amsmath, amssymb, amsfonts, amsthm, mathrsfs, appendix, xcolor, textcomp; the wrapper is the lane's pinned \texttt{amsart} editorial wrapper at 10pt on A4, which restates the theorem-like environments the retained text needs and adds the article's own macro definitions that the retained text needs (\texttt{\textbackslash C}, \texttt{\textbackslash R}, \texttt{\textbackslash Z}), with extra wrapper packages (bm, bbm, dsfont, xspace, cancel, mathtools). The delivered numbering is the unit's own. Assets: no retained span contains a figure environment, a graphics-inclusion command, a PDF-page inclusion, a table or a TikZ picture; no image file is copied into this package, the package declares no asset dependency and no image file is redistributed with it. Licence: version arXiv:2507.22468v1 of this article is distributed under the Creative Commons Attribution 4.0 International licence, as recorded in the arXiv licence metadata for this exact version and shown on the abstract page https://arxiv.org/abs/2507.22468v1. A full rights notice is delivered at licenses/arxiv-2507.22468-CC-BY-4.0.txt. Characters: every retained excerpt is pure ASCII, so the delivered engine, XeLaTeX with its default Latin Modern text font, reports no missing character.
	\begin{align}\label{sampling.def.eq1}
		A\|f\|_{L^p(\R)}^p\le \sum_{n\in \Z} |f(x_n)|^p \le B\|f\|_{L^p(\R)}^p, \ f\in V^p(\phi), 
	\end{align}

	\begin{proposition}\label{prop2}
	A complex-valued vector $x=(x_1, x_2, x_3)\in \C^3$ can be determined, up to a unimodular constant and conjugation, from absolute magnitudes $ |x_n|, n\in \{1,2,3\} $ and relative magnitudes $  |x_n-x_m|, (n,m)\in \{1,2,3\}^2 $.
	\end{proposition}

	\begin{align}\label{eq91}
		f(x_{n+1})\overline{f(x_n)} \ne f(x_n) \overline{f(x_{n+1})}
	\end{align}

\begin{theorem}\label{cpr.adj.thm}
	Let $f\in V^p(\phi)$ be a complex-valued function,  $X=\{x_n\}_{n\in \Z}$ be a set of sampling for $V^p(\phi)$ satisfying \eqref{sampling.def.eq1}  and \eqref{eq91}. Then $f$ can be determined, up to a unimodular constant and conjugation, from absolute magnitude samples $ |f(x_n)|, n\in \Z $ and relative magnitude samples $|f(x_{n+1})-f(x_n)|, |f(x_{n+2})-f(x_n)|, n\in \Z$.
\end{theorem}

\begin{proof}
	Let $g\in V^p(\phi)$ satisfy 
	\begin{align}\label{eq24}
		|f(x_n)| = |g(x_n)|, 
	\end{align}
	\begin{align*}
		|f(x_{n+1}) - f(x_n)| = |g(x_{n+1}) - g(x_n)|
	\end{align*}
	and
	\begin{align*}
		|f(x_{n+2}) - f(x_n)| = |g(x_{n+2}) - g(x_n)|
	\end{align*}
	for all $n\in \Z$.  
	Then
	\begin{align}\label{eq28}
		\Re f(x_{n+1})   \Re f(x_{n})    +  \Im f(x_{n+1})  \Im f(x_{n}) 
		=  \Re g(x_{n+1})   \Re g(x_{n}  )  +  \Im g(x_{n+1})   \Im g(x_{n}) 
	\end{align}
	and
	\begin{align}\label{eq29}
		\Re f(x_{n+2})   \Re f(x_{n})    +  \Im f(x_{n+2})  \Im f(x_{n}) 
		=  \Re g(x_{n+2})   \Re g(x_{n}  )  +  \Im g(x_{n+2})   \Im g(x_{n}) 
	\end{align}
	for all $n\in \Z$.
	By \eqref{eq91}, we have
	\begin{align}\label{eq73}
		\Re f(x_{n+1}) \Im f(x_{n}) \ne \Re f(x_{n}) \Im f(x_{n+1})
	\end{align}
	for all $n\in \Z$.
	Then we know that $f(x_n) \ne 0$ for all $n\in \Z$. Without loss of generality, we assume $f(x_0)=g(x_0)>0$, otherwise replacing $f$ by $e^{i\alpha_1}f$ and $g$ by $e^{i\alpha_2}g$ respectively, where $\alpha_1, \alpha_2\in \R$.
	Obviously, $\Re f(x_0) = \Re g(x_0) > 0$ and $ \Im f(x_0) = \Im g(x_0) =0$. By \eqref{eq24}, \eqref{eq28} and \eqref{eq73}, we have $\Re f(x_1) =  \Re g(x_1)$ and $ \Im f(x_1)= \tau_1 \Im g(x_1) \ne 0$ for some $\tau_1= \pm 1$. Following the similar arguments in part (b) of the proof of Proposition \ref{prop2}, we have
	\begin{align*}
		\Re\big(f(x_{0}), f(x_{1}), f(x_{2})\big) = \Re\big (g(x_{0}), g(x_{1}), g(x_{2})\big)
	\end{align*}
	and
	\begin{align*}
		\Im \big(f(x_{0}), f(x_{1}), f(x_{2}))= \tau_1 \Im (g(x_{0}), g(x_{1}), g(x_{2})\big).
	\end{align*}
	Thus, $(f(x_0), f(x_1), f(x_2))$ is determined, up to a conjugation.
	
	Next we will prove that $(f(x_n))_{n\in \Z}$ is determined, up to a conjugation, 
	%$, i.e.,
	%\begin{align}\label{eq36}
	%	\Re(f(x_n)) = \Re g(x_n), \Im (f(x_n)) = \tau_1 \Im  g(x_n),  n\in \Z
	%\end{align}
	by induction.
	Inductively, we assume
	\begin{align*}
		\Re f(x_n) = \Re g(x_n), \Im f(x_n) = \tau_1 \Im  g(x_n),  n\in [0,n_1]\cap \Z,
	\end{align*}
	where $n_1\ge 2$.
	By \eqref{eq28} and \eqref{eq29},  we have
	\begin{equation*}
		\left\{ \hskip .2in \begin{array}{ll}
			\hskip-.2in
			\big( \Re f(  x_{n_1+1}  )  -  \Re g(  x_{n_1+1}  ) \big)  \Re f(x_{n_1}) + \big(  \Im f(  x_{n_1+1}  )  -  \tau_1 \Im g(  x_{n_1+1}) \big)  \Im f(x_{n_1}  ) {}& \hskip-.10in =  0 \\
			\hskip-.2in
			\big( \Re f(  x_{n_1+1}  )  -  \Re g(  x_{n_1+1}  ) \big)  \Re f(x_{n_1-1}  ) 
			+   \big( \Im  f(  x_{n_1+1}  )  -  \tau_1 \Im g(  x_{n_1+1}  )  \big)  \Im  f(  x_{n_1-1}  ) {}& \hskip-.10in =   0	
		\end{array}\right.
	\end{equation*}
	Together with (\ref{eq73}), we have $\Re f(x_{n_1+1}) = \Re g(x_{n_1+1}) $ and $\Im  f(x_{n_1+1}) = \tau_1 \Im g(x_{n_1+1})$. 
	Therefore, we have 
	\begin{align}\label{cpr.eq.1a}
		\Re f(x_n) = \Re g(x_n) \ \ {\rm and} \ \ \Im f(x_n) = \tau_1 \Im  g(x_n)
	\end{align}
	for all $n\in [0,\infty)\cap \Z$. Following the similar arguments above, we  have \eqref{cpr.eq.1a}  holds for all $n\in (-\infty, 0)\cap \Z$. Therefore 
	\begin{align*}
		f(x_n)=g(x_n) \text{ for all } x_n\in X  \ \ \text{ or }  \ \ f(x_n)=\overline{g(x_n)} \text{ for all } x_n\in X.
	\end{align*}
	As $X$ is a set of sampling for $V^p(\phi)$, we obtain
	\begin{align*}
		f=g \text{ or } f=\overline{g},
	\end{align*}
	which completes the proof.
\end{proof}

\end{document}
